\phantomsection\label{sec:teaching-experience}
\section*{Teaching Experience}
My teaching spans course teaching and student supervision, guided by problem-based learning.
\teachingportfolionote

\subsection*{Course Teaching}
My course teaching follows three core principles: a solid understanding of the subject matter, awareness of students' established experiences and background, and problem-based learning style applied throughout lectures, exercises, and exams. This translates to practices including tailored blog-style literature, inspirational lectures that connect theory to real-world problems, and progressively arranged hands-on exercises.

\pubentry{Aalborg University, 2025 -- 2026}{\href{https://blog.yanlincs.com/ai-system/}{Literature}}
{AI Systems \& Infrastructure}
{}
{Course on streamlined interaction with AI models and the implementation and deployment of scalable, production-ready AI systems on real-world infrastructures.}

\subsection*{Student Supervision}
I am building my supervision experience, supervising Bachelor's and Master's semester projects. My supervision focuses on problem definition, critical thinking about design choices, and delivering convincing conclusions.

\pubentry{Aalborg University, 2026}{\href{https://github.com/Logan-Lin/traj-rep-generalization/releases/latest/download/main.pdf}{Proposal}}
{Trajectory Representation Learning that Generalizes Across Regions, Modes, and Tasks}
{}
{A Master's semester project on learning trajectory representations that generalize across tasks, geographic regions, and transportation modes, modeling vehicle and vessel trajectories relative to road networks and sea depth contours and handling their widely varying sampling intervals.}

\pubentry{Aalborg University, 2026}{\href{https://github.com/Logan-Lin/assess-data-traj/releases/latest/download/main.pdf}{Proposal}}
{Assessment of the Data Foundation for Trajectory Machine Learning}
{}
{A Master's semester project on bringing a data-centric perspective to trajectory machine learning, assessing the quality of training data and developing filtering and augmentation methods to improve model performance on downstream tasks.}

\pubentry{Aalborg University, 2025 -- 2026}{\href{https://github.com/Logan-Lin/camera-sees-behind/releases/latest/download/proposal.pdf}{Proposal}}
{A Camera That Sees Behind}
{}
{A Bachelor's semester project on developing a computationally lightweight model deployable on embedded devices to remove humans from images or video feeds, serving as a prototype GDPR-compliant surveillance system.}
